% ============================================================
% BAB VII - PENUTUP
% Template (ITB STI TA + Proposal.tex):
%   - VII.1 Kesimpulan : empat poin K1..K4, satu lawan satu dengan T-1..T-4
%   - VII.2 Saran      : tindak lanjut + arah pengembangan platform berikutnya
% Rantai keterhubungan: T-N <-> K-N ; tidak menambah klaim baru di luar bab sebelumnya
% ============================================================
\cleardoublepage
\chapter{PENUTUP}
\label{chap:penutup}

Bab ini merangkum capaian penelitian dan menutup rantai keterhubungan yang dijaga sejak rumusan masalah. Setiap butir kesimpulan menjawab tepat satu tujuan T-1 sampai T-4 dengan urutan yang sama, sehingga setiap klaim dapat ditelusuri kembali ke tujuan, kebutuhan, dan skenario uji yang bersesuaian. Kesimpulan disusun pada tingkat capaian yang telah terverifikasi secara fungsional dan struktural pada lingkungan node tunggal, sedangkan karakterisasi beban kuantitatif pada klaster multi-node ditempatkan sebagai arah pengembangan sesuai lingkup batasan B-5 pada Subbab~\ref{sec:batasan}. Saran kemudian mengarahkan langkah lanjutan untuk memperluas karakterisasi beban dan untuk mengembangkan platform pada tahap berikutnya.

\section{Kesimpulan}
\label{sec:kesimpulan}

Penelitian ini menghasilkan satu instansiasi arsitektur DataOps dan MLOps terintegrasi pada Kubernetes yang dievaluasi melalui \textit{use case} pasar aset kripto sebagai instans verifikasi. Berdasarkan hasil evaluasi pada Bab~\ref{chap:evaluasi}, empat tujuan platform dapat disimpulkan sebagai berikut.

\begin{enumerate}
  \item \textbf{Arsitektur terintegrasi dan GitOps (T-1)} \\
  Seluruh komponen platform terbukti dapat dipasang dan direkonsiliasi dari satu sumber Git oleh Argo CD, dengan perubahan definisi konvergen tanpa \textit{drift} dan rekonstruksi penuh pada \textit{node} k3s baru tanpa intervensi manual selain \textit{bootstrap} awal. Capaian ini menjawab RM-1 dengan menggantikan fragmentasi tumpukan teknologi oleh satu bidang kendali deklaratif yang dapat diaudit.
  \item \textbf{Tata kelola data otomatis (T-2)} \\
  Katalog metadata, \textit{lineage}, dan kontrol kualitas terbukti berjalan sebagai layanan bersama sepanjang \textit{pipeline} data, fitur, dan model, dengan peristiwa \textit{lineage} OpenLineage dan kontrak Great Expectations tertanam pada jalur produksi. Capaian ini menjawab RM-2 dengan menjadikan tata kelola bagian melekat dari \textit{pipeline}, bukan pekerjaan pasca-fakta.
  \item \textbf{Deteksi \textit{drift} dan \textit{point-in-time correctness} (T-3)} \\
  \textit{Control loop} pelatihan ulang terbukti terpicu otomatis ketika injeksi \textit{drift} terkendali diterapkan pada satu fitur dan penyusunan fitur historis terbukti menolak kebocoran temporal pada jalur pelatihan. Capaian ini menjawab RM-3, sedangkan pengamatan ketepatan deteksi pada pergeseran distribusi alami berdurasi panjang menjadi arah pengembangan lanjutan.
  \item \textbf{Layanan fitur \textit{dual-store} (T-4)} \\
  Jalur penyajian fitur \textit{online}, \textit{offline}, dan vektor terbukti terbentuk dengan kontrak fitur yang konsisten antara fase pelatihan dan fase \textit{inference}. Capaian ini menjawab RM-4 pada tingkat fungsional dan struktural, sedangkan karakterisasi latensi p99 dan \textit{throughput} pada beban tinggi di klaster multi-node menjadi arah pengembangan lanjutan.
\end{enumerate}

Keempat kesimpulan tersebut menegaskan bahwa kontribusi utama penelitian berupa instansiasi arsitektur terintegrasi telah tercapai dan tertelusur secara penuh dari rumusan masalah hingga skenario uji. Pemenuhan yang dilaporkan bersifat fungsional dan struktural pada lingkungan node tunggal, sehingga setiap klaim disusun agar tepat sepadan dengan bukti yang telah diamati tanpa melampauinya. Nilai utama instansiasi ini terletak pada pemindahan beban integrasi yang sebelumnya berulang menjadi satu deskripsi deklaratif tunggal, sehingga rekonstruksi lapis kendali tidak lagi bergantung pada pengetahuan tersembunyi satu pengembang. Platform terbukti dapat direkonstruksi, dikelola secara deklaratif, dan menjalankan siklus hidup data serta model secara terotomasi, sedangkan penilaian kinerja menyeluruh menjadi pekerjaan lanjutan yang terdefinisi jelas.

\section{Saran}
\label{sec:saran}

Saran pertama diarahkan pada perluasan karakterisasi beban kuantitatif platform. Pengujian beban penuh perlu dijalankan pada lingkungan multi-node untuk memverifikasi skalabilitas horizontal, latensi p99, dan \textit{throughput} puncak yang pada penelitian ini terverifikasi pada tingkat fungsional dan struktural di lingkungan node tunggal. Pengamatan \textit{drift} juga perlu diperpanjang durasinya pada aliran data nyata agar ketepatan deteksi terhadap pergeseran distribusi alami berdurasi panjang dapat dinilai melengkapi verifikasi melalui injeksi \textit{drift} terkendali. Selain itu, penilaian kebenaran nilai fitur sebaiknya dilakukan setelah jalur ingestasi sampai lapis \textit{gold} terisi penuh sehingga seluruh kontrak fitur dapat diuji pada kondisi data lengkap.

Saran kedua diarahkan pada arah pengembangan platform berikutnya di luar lingkup penelitian ini. Beberapa arah yang teridentifikasi mencakup perluasan dukungan multi-\textit{tenancy} pada lapis tata kelola, otomasi kebijakan kuota pada Kueue untuk beban pelatihan yang lebih besar, integrasi metrik biaya OpenCost dengan kebijakan \textit{retraining}, serta adopsi profil keamanan Kyverno yang lebih ketat setelah profil dasar matang. Arah tersebut melanjutkan sifat \textit{domain-agnostic} platform sehingga \textit{use case} lain dapat diadopsi tanpa mengubah lapis kendali. Platform dapat berkembang dari instansiasi yang terverifikasi struktural menjadi fondasi yang teruji pada beban produksi nyata.
